\section*{Chapitre \ref{ch:iv1:continuous-and-geometric-probability} --- Probabilités continues et géométriques}
\markboth{\protect\omTocMark{}{Réponses~: Probabilités continues et géométriques}}{\protect\omTocMark{}{Réponses~: Probabilités continues et géométriques}}

\begin{solution}{iq:iv1:continuous-and-geometric-probability:1}
La plus longue des deux parties a une espérance de $3/4$. Le rapport des espérances de la plus courte et de la plus longue vaut $2\ln 2 - 1 \approx 0.386$.

La plus longue partie suit une loi uniforme sur $[\tfrac12, 1]$, de densité 2, donc son espérance vaut $3/4$. Si $\ell$ désigne la longueur de la plus longue partie, le rapport vaut $(1 - \ell)/\ell$, et
\[
E\Bigl[\frac{1 - \ell}{\ell}\Bigr] = \int_{1/2}^{1} 2\Bigl(\frac1\ell - 1\Bigr) d\ell = 2\ln 2 - 1.
\]
\iqwrong{Diviser les deux espérances, soit $(1/4)/(3/4) = 1/3$. L'espérance d'un rapport n'est pas le rapport des espérances, car un rapport varie davantage lorsque le dénominateur est petit.}
\end{solution}

\begin{solution}{iq:iv1:continuous-and-geometric-probability:2}
La probabilité vaut $\tfrac12 + \tfrac12 \ln 2 \approx 0.847$.

Traçons le carré unité et la courbe $y = 1/(2x)$. Pour $x \le \tfrac12$, la courbe se situe au-dessus du carré, donc tout $y$ convient, ce qui apporte $\tfrac12$. Pour $x > \tfrac12$, la hauteur admissible vaut $1/(2x)$, ce qui apporte
\[
\int_{1/2}^{1} \frac{dx}{2x} = \tfrac12 \ln 2.
\]
\iqwrong{Répondre $1/2$, en traitant $XY$ comme s'il était uniforme. Le produit est biaisé vers le bas~: sa médiane vaut seulement environ $0.19$, de sorte que $1/2$ se situe bien dans sa queue supérieure.}
\end{solution}

\begin{solution}{iq:iv1:continuous-and-geometric-probability:3}
L'espérance vaut $1/3$ et la variance vaut $1/18$.

L'événement $|X - Y| > t$ correspond à deux triangles d'angle, d'aire totale $(1 - t)^2$. La densité de $D = |X - Y|$ vaut donc $2(1 - t)$ sur $[0, 1]$, et $E[D] = \int_0^1 2t(1 - t)\,dt = 1/3$.

Pour la variance, $E[D^2] = E[(X - Y)^2] = 2\,\mathrm{Var}(X) = 1/6$, d'où une variance égale à $1/6 - 1/9 = 1/18$.

\iqfollow{Et pour deux normales centrées réduites indépendantes~?}{$X - Y$ a pour variance 2, donc $E|X - Y| = \sqrt{2}\sqrt{2/\pi} = 2/\sqrt{\pi}$.}
\end{solution}

\begin{solution}{iq:iv1:continuous-and-geometric-probability:4}
La probabilité vaut $7/16$.

Le couple des instants d'arrivée suit une loi uniforme sur un carré de côté 1 heure\index{probabilité géométrique}. Ils se rencontrent lorsque $|X - Y| \le \tfrac14$, soit une bande autour de la diagonale. En dehors de la bande se trouvent deux triangles d'angle, chacun de côtés de l'angle droit $\tfrac34$, dont l'aire totale vaut $(\tfrac34)^2 = 9/16$.

S'ils attendent chacun toute l'heure, la formule donne 1, et s'ils n'attendent pas du tout, elle donne 0, ce qui constitue une vérification rapide.

\iqwrong{La moitié, en doublant une fenêtre de 15 minutes pour obtenir 30 minutes sur 60. La fenêtre est tronquée aux extrémités de l'heure, pour quelqu'un qui arrive près de midi ou de 13~h.}
\end{solution}

\begin{solution}{iq:iv1:continuous-and-geometric-probability:5}
À moins de $0.05$ de la frontière~: $0.1$, environ $0.651$ et environ $0.99997$. Dans la boule~: environ $0.524$ pour $d = 3$, et environ $0.0025$ pour $d = 10$.

Le point est à au moins $0.05$ de la frontière seulement si chaque coordonnée appartient à $[0.05, 0.95]$, ce qui a pour probabilité $0.9$ par coordonnée, donc la probabilité vaut $0.9^d$. La boule de diamètre 1 a pour volume $\pi/6$ pour $d = 3$, et pour $d = 10$ il vaut $\pi^5/(5! \cdot 2^{10})$.

En grande dimension, presque tout le cube se trouve donc près de ses sommets et de ses arêtes\index{fléau de la dimension}, et presque rien près du centre. Les points de données sont éloignés les uns des autres et aussi de tout centre~: la distance typique au centre du cube unité croît comme $\sqrt{d/12}$, soit environ $2.9$ pour $d = 100$.

\iqwhy{C'est pourquoi les méthodes des plus proches voisins et les estimations de densité échouent lorsqu'il y a beaucoup de variables~: avec 100 variables, il n'existe pas de point «~voisin~», et chaque point ressemble à une valeur aberrante.}
\end{solution}

\begin{solution}{iq:iv1:continuous-and-geometric-probability:6}
Une double cassure donne $1/4$. Casser la partie la plus longue donne $2\ln 2 - 1 \approx 0.386$, et casser une partie choisie au hasard donne la moitié de cette valeur, soit environ $0.193$.

Trois morceaux forment un triangle exactement lorsque chacun est inférieur à $\tfrac12$, puisque le plus long doit être plus petit que la somme des deux autres, qui vaut $1$ moins lui-même. Pour deux coupures en $u < v$, cela exige $u < \tfrac12$, $v - u < \tfrac12$ et $v > \tfrac12$. Dans le carré unité, il s'agit de la zone hachurée de la \cref{fig:iv1:stick}, dont l'aire vaut $1/4$.

Pour la partie la plus longue, notons sa longueur $\ell$. La partie la plus courte est déjà inférieure à $\tfrac12$. Une coupure en $a$ laisse deux morceaux inférieurs à $\tfrac12$ lorsque $\ell - \tfrac12 < a < \tfrac12$, ce qui représente une fraction $(1 - \ell)/\ell$ de la partie. En moyennant sur $\ell$, on retrouve la même intégrale que pour le rapport des morceaux~: $2\ln 2 - 1$.

Choisir une partie au hasard retient la plus longue dans la moitié des cas. Casser la plus courte ne peut jamais réussir, car la partie la plus longue, laissée intacte, dépasse déjà $\tfrac12$. La probabilité vaut donc la moitié de la valeur précédente.

\begin{omfigure}
\begin{tikzpicture}[x=3.4cm,y=3.4cm]
\fill[omDef!25] (0,0.5) -- (0.5,1) -- (0.5,0.5) -- cycle;
\fill[omDef!25] (0.5,0) -- (1,0.5) -- (0.5,0.5) -- cycle;
\draw (0,0) rectangle (1,1);
\draw[gray,dashed] (0.5,0) -- (0.5,1);
\draw[gray,dashed] (0,0.5) -- (1,0.5);
\draw[omDef] (0,0.5) -- (0.5,1);
\draw[omDef] (0.5,0) -- (1,0.5);
\node[left,font=\footnotesize] at (0,0.5) {$\tfrac12$};
\node[left,font=\footnotesize] at (0,1) {$1$};
\node[below,font=\footnotesize] at (0.5,0) {$\tfrac12$};
\node[below,font=\footnotesize] at (1,0) {$1$};
\node[below,font=\footnotesize] at (0,0) {$0$};
\node[below=5.5mm,font=\footnotesize] at (0.25,0) {première coupure};
\node[left=3mm,font=\footnotesize,rotate=90,anchor=south] at (0,0.25) {seconde coupure};
\end{tikzpicture}

\omcaption{Deux coupures indépendantes sur un bâton. Les morceaux forment un triangle dans les deux triangles hachurés, soit un quart du carré.}\label{fig:iv1:stick}
\end{omfigure}

\iqwrong{Traiter le second problème comme le premier. Choisir la partie la plus longue rend le triangle plus probable, car cela écarte le pire cas.}
\end{solution}

\begin{solution}{iq:iv1:continuous-and-geometric-probability:7}
La probabilité vaut $1/2$.

Ils ne forment pas un triangle exactement lorsque l'un est plus long que la somme des deux autres. Deux de ces trois événements ne peuvent jamais se produire ensemble, puisque seul le bâton le plus long peut violer la règle. Chaque événement, tel que $Z > X + Y$, a pour probabilité $1/6$~: il s'agit du coin du cube unité en $(0, 0, 1)$, une pyramide de volume $\tfrac13 \cdot \tfrac12 \cdot 1$. La probabilité d'échec vaut donc $3 \cdot \tfrac16 = \tfrac12$.

\iqwrong{Répondre $1/4$, en se rappelant le bâton brisé. Dans ce problème, les trois longueurs doivent totaliser 1, et elles ne sont pas indépendantes. Ici, chacune peut prendre n'importe quelle valeur, et la probabilité de former un triangle est plus grande.}
\iqfollow{Avec quatre bâtons et une figure à quatre côtés~?}{Le plus long doit être inférieur à la somme des trois autres. Chaque événement d'échec a pour probabilité $1/24$, donc la réponse vaut $1 - 4/24 = 5/6$.}
\end{solution}

\begin{solution}{iq:iv1:continuous-and-geometric-probability:8}
Une aiguille de longueur 1 coupe une ligne avec une probabilité de $2/\pi \approx 0.637$. Une aiguille de longueur 3 réalise en moyenne $6/\pi \approx 1.91$ intersections.

Pour la première, notons $\theta$ l'angle avec les lignes, uniforme sur $[0, \pi)$. L'étendue de l'aiguille perpendiculairement aux lignes vaut $\sin\theta$, et pour cet angle la probabilité de coupure vaut $\sin\theta$, car la position de l'aiguille par rapport aux lignes est uniforme\index{aiguille de Buffon}. Sa moyenne vaut $2/\pi$.

Pour la seconde, utilisons la linéarité\index{linéarité de l'espérance}\index{nouille de Buffon}. Une aiguille de longueur 3 équivaut à trois aiguilles de longueur 1 mises bout à bout. Chacune a le même nombre moyen d'intersections, quelles que soient les jonctions, donc la moyenne vaut $3 \cdot 2/\pi$. En fait, cela vaut pour toute forme~: le nombre moyen d'intersections vaut $2L/\pi d$, où $L$ est la longueur et $d$ l'écart entre les lignes. Un cercle de diamètre 1 a $L = \pi$ et coupe les lignes deux fois à coup sûr, et la formule donne 2, ce qui constitue une belle démonstration de la constante.

\iqwrong{Chercher $P(\text{au moins une})$ pour l'aiguille longue. Cela demande une intégrale sur les angles, alors que la moyenne n'en demande pas. La question porte sur la moyenne parce que la linéarité la rend presque gratuite.}
\end{solution}

\begin{solution}{iq:iv1:continuous-and-geometric-probability:9}
La probabilité vaut $n/2^{n - 1}$, soit $3/4$ pour $n = 3$ et $5/16$ pour $n = 5$. Un triangle construit sur trois points contient le centre avec une probabilité de $1/4$.

Si les points se trouvent dans un demi-cercle, alors en parcourant le cercle dans le sens horaire il existe un premier point unique de l'arc\index{argument du demi-cercle}. Les événements «~le point $i$ est le premier~» sont disjoints. Sachant que le point $i$ est le premier, tous les autres doivent tomber dans le demi-cercle situé dans le sens horaire à partir de lui, ce qui a pour probabilité $2^{-(n - 1)}$. Il existe $n$ tels événements, donc le total vaut $n/2^{n - 1}$.

Un triangle contient le centre exactement lorsque ses trois sommets ne tiennent pas dans un demi-cercle, soit $1 - 3/4 = 1/4$.

\iqwrong{Fixer le premier point et exiger que les autres se trouvent dans le demi-cercle qui le suit, ce qui donne $2^{-(n - 1)}$. On oublie ainsi que n'importe lequel des $n$ points peut être le début de l'arc.}
\iqfollow{Et pour $n = 2$~?}{La formule donne 1, ce qui est juste~: deux points se trouvent toujours dans un demi-cercle.}
\end{solution}

\begin{solution}{iq:iv1:continuous-and-geometric-probability:10}
Deux cordes se croisent avec une probabilité de $1/3$, et 10 cordes présentent en moyenne 15 croisements.

Numérotons les quatre points dans l'ordre autour du cercle. Quelles que soient leurs positions, trois appariements sont possibles, chacun étant également probable~: le premier point est apparié à son voisin dans le sens horaire, à celui qui lui est opposé, ou à son voisin dans le sens antihoraire. Seul l'appariement opposé se croise\index{argument de symétrie}, donc la chance vaut $1/3$.

Il existe $\binom{10}{2} = 45$ paires de cordes, et chaque paire correspond à une configuration de quatre points aléatoires indépendants. Par linéarité\index{linéarité de l'espérance}, quelle que soit la dépendance entre les paires, le nombre moyen de croisements vaut donc $45 \cdot \tfrac13 = 15$.

\iqwrong{Tenter un calcul sur les longueurs des arcs. Seul l'ordre relatif des quatre points compte, ce qui explique que la réponse ne comporte aucune intégrale.}
\iqfollow{Les croisements sont-ils indépendants~?}{Non, mais la linéarité n'en a pas besoin. Une question sur la variance demanderait les corrélations, et là le travail commence.}
\end{solution}

\begin{solution}{iq:iv1:continuous-and-geometric-probability:11}
Les réponses sont $1/3$, $1/2$ et $1/4$, selon la méthode. C'est le paradoxe de Bertrand\index{paradoxe de Bertrand}.

Le côté du triangle équilatéral inscrit mesure $\sqrt{3}$, et une telle corde se trouve à une distance de $\tfrac12$ du centre. La corde est donc plus longue exactement lorsqu'elle passe à moins de $\tfrac12$ du centre.

Extrémités aléatoires~: on fixe une extrémité. L'autre doit se trouver sur le tiers du cercle opposé à celle-ci~: $1/3$.

Distance aléatoire le long d'un rayon~: la distance au centre est uniforme sur $[0, 1]$, donc la probabilité d'être inférieure à $\tfrac12$ vaut $1/2$.

Milieu aléatoire dans le disque~: le milieu doit se trouver dans le disque de rayon $\tfrac12$, qui a un quart de l'aire~: $1/4$.

Les trois sont des réponses justes à des questions différentes. «~Au hasard~» ne définit pas une loi de probabilité sur un ensemble infini. La première chose à faire lorsqu'une question emploie cette expression est de se demander quel mécanisme a produit l'objet.

\iqwhy{Le même problème se pose chaque fois que l'on attribue un a priori plat. Plat en volatilité et plat en variance sont des croyances différentes, et le choix de l'une ou l'autre peut modifier un résultat. Énoncer le paramétrage avant d'énoncer la réponse.}
\end{solution}

\begin{solution}{iq:iv1:continuous-and-geometric-probability:12}
La première probabilité vaut $1/12$, et la seconde $3/4$.

Deux normales standard indépendantes\index{symétrie rotationnelle des normales} ont une densité qui ne dépend que de la distance à l'origine, donc la direction de $(X, Y)$ est uniforme. Chaque événement devient un secteur, dont la probabilité est son angle rapporté à l'angle total.

Pour $X > 0$ et $Y > \sqrt{3}X$~: la droite $y = \sqrt{3}x$ fait un angle de $60$ degrés, et la région s'étend de là jusqu'à l'axe vertical, soit un secteur de $30$ degrés. La chance vaut donc $30/360 = 1/12$.

Pour $X + Y > 0$ sachant $X > 0$~: le demi-plan $x > 0$ correspond à $180$ degrés. Le demi-plan $x + y > 0$ le recoupe dans le secteur allant de $-45$ à $90$ degrés, soit $135$ degrés. La chance vaut donc $135/180 = 3/4$.

\iqwrong{Conditionner sur $X$ et intégrer une queue de loi normale. Cela fonctionne, mais c'est lent, et l'on manque ce qui rend les lois normales particulières.}
\iqfollow{En termes de corrélation~?}{$X + Y$ et $X$ ont une corrélation de $1/\sqrt{2}$, et la formule de concordance des signes $\tfrac12 + \arcsin(\rho)/\pi$ donne le même $3/4$.}
\end{solution}

\begin{solution}{iq:iv1:continuous-and-geometric-probability:13}
$P(X^2 + Y^2 > 4) = e^{-2} \approx 0.135$. Le plus grand a pour espérance $1/\sqrt{\pi} \approx 0.564$ et pour variance $1 - 1/\pi \approx 0.682$.

En coordonnées polaires, $R^2 = X^2 + Y^2$ suit une loi exponentielle d'espérance 2\index{coordonnées polaires pour les normales}, donc la probabilité qu'il dépasse 4 vaut $e^{-4/2}$.

Écrivons le plus grand $\max(X, Y) = \tfrac12(X + Y) + \tfrac12|X - Y|$. La première partie est d'espérance nulle, et $X - Y$ suit une loi normale de variance 2, donc $E|X - Y| = \sqrt{2}\sqrt{2/\pi} = 2/\sqrt{\pi}$. L'espérance vaut donc $1/\sqrt{\pi}$.

Pour le second moment, une symétrie d'une ligne suffit\index{argument de symétrie}. En remplaçant $(X, Y)$ par $(-X, -Y)$, le plus grand devient l'opposé du plus petit, donc les carrés du plus grand et du plus petit ont la même espérance. Leur somme vaut $X^2 + Y^2$, d'espérance 2, donc chacun a pour espérance 1. La variance vaut $1 - 1/\pi$.

\iqwhy{Le maximum et le minimum de rendements gaussiens interviennent dans le risque de queue et les options best-of. Une option best-of sur deux rendements gaussiens indépendants et centrés vaut environ $0.56$ écart type.}
\end{solution}

\begin{solution}{iq:iv1:continuous-and-geometric-probability:14}
La machine de taux 1 tombe en panne en dernier avec une probabilité de $7/12$. Les trois ont cessé de fonctionner après $73/60 \approx 1.22$ année en moyenne.

Utilisons la propriété de course des lois exponentielles\index{exponentielles en compétition}~: la prochaine panne parmi celles encore en fonctionnement est celle de la machine $i$ avec une probabilité proportionnelle à son taux, et après une panne les autres repartent de zéro, par absence de mémoire.

Pour que la machine de taux 1 soit la dernière, soit 2 tombe en premier (chance $2/6$) puis 3 bat 1 (chance $3/4$), soit 3 tombe en premier (chance $3/6$) puis 2 bat 1 (chance $2/3$). Cela fait $\tfrac14 + \tfrac13 = \tfrac{7}{12}$.

Pour l'espérance du maximum, utilisons l'inclusion--exclusion\index{inclusion--exclusion} sur les durées. Le minimum d'un ensemble suit une loi exponentielle de taux égal à la somme des taux, donc
\[
E[\max] = 1 + \tfrac12 + \tfrac13 - \tfrac13 - \tfrac14 - \tfrac15 + \tfrac16 = \tfrac{73}{60}.
\]
\iqwrong{Prendre l'espérance de la machine la plus lente, soit 1 an, ce qui oublie que les autres peuvent tomber en panne plus tard, ou additionner toutes les espérances. Le maximum est inférieur à la somme des espérances, et supérieur à chacune d'elles.}
\end{solution}

\begin{solution}{iq:iv1:continuous-and-geometric-probability:15}
La probabilité d'en avoir au moins une vaut $1 - e^{-1} \approx 0.632$. Sachant qu'il y en a exactement 3 dans la minute, la probabilité qu'aucune ne soit arrivée dans la première moitié vaut $1/8$, et la première est arrivée après 15 secondes en moyenne.

En 30 secondes à raison de 2 par minute, le nombre suit une loi de Poisson de moyenne 1, donc la chance qu'il n'y en ait aucune vaut $e^{-1}$.

Sachant le nombre\index{processus de Poisson}, les instants d'arrivée sont indépendants et uniformes sur la minute, quel que soit le taux. Les 3 tombent donc dans la seconde moitié avec une chance de $(1/2)^3$. La première est le minimum de 3 uniformes, d'espérance égale à $1/4$ de minute. Une simulation des écarts exponentiels réels, en ne gardant que les minutes comportant exactement 3 arrivées, le confirme.

\iqwrong{Utiliser l'écart moyen de 30 secondes donné par le taux. Une fois qu'on sait qu'il y en a trois dans la minute, les écarts sont plus courts, et le taux n'apporte plus rien.}
\iqwhy{Cela permet de tester si des événements suivent réellement un processus de Poisson~: sachant le nombre, les instants doivent sembler uniformes. Un regroupement autour des annonces apparaît aussitôt.}
\end{solution}

\begin{solution}{iq:iv1:continuous-and-geometric-probability:16}
Le rapport du plus petit au plus grand est uniforme sur $[0, 1]$. L'espérance de $X/Y$ est infinie, ce qui s'accorde très bien avec une médiane de 1.

Pour $t \in (0, 1)$, l'événement $\min/\max < t$ correspond à l'aire des deux triangles proches des axes\index{probabilité géométrique}~: sous la droite $y = tx$ et au-dessus de la droite $y = x/t$. Chacun a pour aire $t/2$, donc la chance vaut $t$. Le rapport est donc uniforme, et son espérance vaut $1/2$.

Pour $X/Y$~: par symétrie, il est aussi probable d'être inférieur à 1 que supérieur à 1, donc la médiane vaut 1. Mais $E[X/Y] = E[X]\,E[1/Y] = \tfrac12 \int_0^1 dy/y = \infty$. La queue de la distribution est trop lourde, car $Y$ peut être très petit.

\iqwrong{Donner $E[X]/E[Y] = 1$ comme espérance du rapport. L'indépendance donne le produit de $E[X]$ et de $E[1/Y]$, et $E[1/Y]$ n'est pas $1/E[Y]$.}
\iqwhy{Moyenner des rapports est dangereux pour cette raison, car un petit dénominateur peut dominer. Moyenner des rendements sur de petites bases de capital, ou un ratio cours-bénéfice moyen sur des entreprises aux bénéfices minuscules, pose le même problème.}
\end{solution}

\begin{solution}{iq:iv1:continuous-and-geometric-probability:17}
L'erreur type vaut environ $1.642/\sqrt{n}$, et pour $\pm 0.001$ à 95~\% il faut environ $1.04 \times 10^7$ points.

Un point tombe dans le quart de disque avec la probabilité $p = \pi/4$. La proportion de succès a pour variance $p(1 - p)/n$. L'estimation vaut 4 fois cette proportion, donc son erreur type vaut $4\sqrt{p(1 - p)/n} = 1.642/\sqrt{n}$.

Pour obtenir $1.96 \cdot 1.642/\sqrt{n} = 0.001$, il faut $n = (1.96 \cdot 1.642/0.001)^2 \approx 1.04 \times 10^7$. Chaque chiffre supplémentaire exige 100 fois plus de points\index{erreur de Monte Carlo}.

\iqwrong{Penser que l'erreur vaut $1/n$, comme pour une somme que l'on additionne, ou qu'un chiffre de plus coûte 10 fois plus de calcul. L'erreur décroît avec la racine carrée de $n$, ce qui rend la méthode de Monte Carlo classique coûteuse.}
\iqfollow{Peut-on faire mieux sans plus de points~?}{Oui. On peut stratifier le carré en une grille, utiliser une variable de contrôle régulière, ou des points à faible discrépance. Chacune donne la même réponse avec une variance plus faible.}
\end{solution}

\begin{solution}{iq:iv1:continuous-and-geometric-probability:18}
La probabilité vaut $1/4$, et la formule générale est $\tfrac18 + \tfrac{3}{4\pi} \arcsin \rho$. Pour $\rho = -\tfrac12$ elle atteint 0, ce qui correspond à la plus faible corrélation possible pour trois variables.

Comptons les désaccords de signe\index{probabilité de l'orthant}. Deux des variables diffèrent en signe avec la probabilité $d = \tfrac12 - \arcsin(\rho)/\pi$, qui dépend de l'angle entre elles. Soit $a$ la chance que les trois soient positives, qui est aussi la chance qu'elles soient toutes négatives, par symétrie. Un triplet de signes mélangés est en désaccord sur exactement 2 de ses 3 paires, donc en sommant sur les paires,
\[
3d = 2(1 - 2a), \qquad\text{ce qui donne}\qquad a = \tfrac12 - \tfrac34 d = \tfrac18 + \frac{3}{4\pi}\arcsin\rho.
\]
Pour $\rho = \tfrac12$, l'arcsinus vaut $\pi/6$, et cela donne $\tfrac18 + \tfrac18$.

Une seconde voie pour cette valeur~: on écrit $X_i = (W_0 - W_i)/\sqrt{2}$ avec des normales standard indépendantes $W_0, \dots, W_3$. Elles ont une corrélation $\tfrac12$, et les trois sont positives exactement lorsque $W_0$ est la plus grande des quatre, ce qui arrive avec probabilité $1/4$.

Contrôles aux limites~: $\rho = 0$ donne $1/8$, et $\rho = 1$ donne $1/2$. Pour $\rho = -\tfrac12$, on a $X + Y + Z = 0$ avec certitude, puisque la variance de la somme vaut $3 + 6\rho = 0$. Elles ne peuvent pas toutes être positives, et une corrélation plus faible rendrait cette variance négative\index{limite d'équicorrélation}. La formule a donc le bon signe exactement à la limite de faisabilité.

\iqwrong{Donner $\tfrac12 \cdot \tfrac23 \cdot \tfrac23 = 2/9$, en enchaînant les probabilités de signe par paires. Sachant que $X$ et $Y$ sont toutes deux positives, $Z$ a plus de chances d'être positive que si l'on ne connaît que $X$~: la probabilité vaut $3/4$, et non $2/3$.}
\iqfollow{Et $n$ variables à corrélation $\tfrac12$~?}{La même construction donne $1/(n + 1)$. Pour un $\rho$ général et $n \ge 4$, il n'existe pas de formule aussi simple, et le résultat à trois variables est une propriété particulière de ce cas.}
\end{solution}

\begin{solution}{iq:iv1:continuous-and-geometric-probability:19}
Le plus court a pour moyenne $1/100$, et le plus long a pour moyenne $H_{10}/10 \approx 0.293$.

On utilise la représentation exponentielle\index{espacements uniformes}~: les $n$ espacements d'un bâton coupé en $n - 1$ points uniformes ont la même loi que $n$ exponentielles indépendantes divisées par leur somme, et cette somme est indépendante de la forme. On peut donc calculer avec des exponentielles, puis diviser par la moyenne de la somme, qui vaut $n$.

Pour des exponentielles de moyenne 1, la plus petite des $n$ a pour moyenne $1/n$. En retirant ce minimum à toutes, les $n - 1$ autres sont de nouvelles exponentielles, par absence de mémoire. Les valeurs triées ont donc des écarts de moyennes $1/n$, $1/(n-1)$, \dots, $1$\index{représentation de Rényi}. La plus grande a donc pour moyenne $H_n$.

On divise par $n$. Le plus petit morceau a pour moyenne $1/n^2 = 1/100$, et le plus long $H_n/n = 2.929/10$. Le plus court vaut un dixième du morceau moyen de $0.1$, et le plus long vaut environ trois fois cette moyenne.

\iqwrong{Donner $1/10$ pour le plus petit, ce qui est la position moyenne de la première coupure. C'est la moyenne d'un morceau donné. Le plus petit est le minimum de dix de ces morceaux, et sa moyenne est beaucoup plus petite.}
\iqwhy{Cela explique l'allure des écarts entre événements lorsqu'ils arrivent au hasard. Avec 10 écarts à des instants aléatoires, on s'attend à ce que l'un d'eux vaille 3 fois la moyenne, et un autre un dixième de la moyenne~: aucun de ces deux faits n'a de sens en soi.}
\end{solution}

\begin{solution}{iq:iv1:continuous-and-geometric-probability:20}
La probabilité vaut $11/16$. Pour $k$ morceaux, elle vaut $1 - k/2^{k - 1}$.

Des morceaux peuvent former un polygone fermé\index{inégalité du polygone} exactement lorsqu'aucun morceau n'est au moins égal à la somme des autres. Comme le total vaut 1, cela revient à dire qu'aucun morceau n'a une longueur au moins égale à $\tfrac12$. Deux morceaux ne peuvent pas valoir tous deux $\tfrac12$ ou plus, sauf sur un ensemble négligeable, donc ces événements sont disjoints.

Un morceau donné vaut au moins $\tfrac12$ avec probabilité $(\tfrac12)^{k - 1}$~: pour le premier morceau, les $k - 1$ coupures doivent toutes se trouver dans la moitié haute, et les morceaux sont échangeables. La probabilité qu'un morceau le fasse vaut donc $k/2^{k - 1}$, et la réponse en est le complémentaire. Pour $k = 5$, cela donne $1 - 5/16$, et pour $k = 3$, le $1/4$ familier.

Le lien avec les points sur un cercle est exact. Couper un cercle en $k$ points uniformes donne $k$ arcs ayant la même loi conjointe que ces morceaux. Les $k$ points se trouvent tous dans un demi-cercle exactement lorsqu'un arc vaut au moins $\tfrac12$. Les deux formules expriment donc un seul et même fait.

\iqwrong{Multiplier le $1/4$ du triangle par un facteur pour chaque morceau supplémentaire, comme si les choses devenaient plus difficiles. Avec davantage de morceaux, c'est plus facile, car il est moins probable qu'un morceau donné soit très long.}
\end{solution}

\begin{solution}{iq:iv1:continuous-and-geometric-probability:21}
L'aire moyenne vaut $\frac{3}{2\pi} \approx 0.477$.

Notons $a$, $b$ et $c$ les arcs entre points voisins, dont la somme vaut $2\pi$. On relie les trois sommets au centre. L'aire du triangle est la somme de trois triangles dont les côtés valent 1 et dont l'angle au sommet est celui d'un arc\index{arcs du triangle inscrit}~:
\[
\text{Aire} = \tfrac12 (\sin a + \sin b + \sin c),
\]
Cette formule reste valable même lorsque le centre est à l'extérieur, car le sinus d'un arc supérieur à $\pi$ est négatif, et la formule retranche alors cette part. Les arcs ont la même loi marginale par symétrie. Un arc seul a pour densité $\frac{1}{\pi}\bigl(1 - \frac{a}{2\pi}\bigr)$ sur $(0, 2\pi)$, donc
\[
E[\sin a] = \int_0^{2\pi} \frac{\sin a}{\pi}\Bigl(1 - \frac{a}{2\pi}\Bigr) da = \frac{1}{\pi}.
\]
L'espérance de l'aire vaut donc $\tfrac32 \cdot \tfrac1\pi$. Une simulation de l'aire à partir des coordonnées concorde avec ce résultat, et elle confirme la formule signée point par point. Le plus grand triangle possible, l'équilatéral, a pour aire $3\sqrt{3}/4 \approx 1.30$, donc la moyenne vaut un peu plus d'un tiers du meilleur.

\iqwrong{Prendre l'aire du triangle dont les arcs valent en moyenne $2\pi/3$ chacun, ce qui donne l'équilatéral $1.30$. L'aire est concave près du centre du simplexe, donc la moyenne reste bien en deçà de cette valeur.}
\end{solution}

\begin{solution}{iq:iv1:continuous-and-geometric-probability:22}
La probabilité vaut $\frac{4\sqrt{2} - 5}{3} \approx 0.219$.

On place le centre à l'origine, de sorte que le carré soit $[-\tfrac12, \tfrac12]^2$. Par symétrie, on considère la tranche du carré la plus proche du côté droit, où $|y| \le x$~; il y a quatre tranches de ce type, et chacune est un triangle allant du centre à un côté. Un point de cette tranche est à distance $\tfrac12 - x$ du bord, et il est plus proche du centre si
\[
\sqrt{x^2 + y^2} < \tfrac12 - x \iff x < \tfrac14 - y^2,
\]
il s'agit d'une parabole\index{parabole comme lieu}, le lieu des points à égale distance d'un point et d'une droite. La courbe rencontre le bord de la tranche $y = x$ là où $x^2 + x - \tfrac14 = 0$, soit en $y_0 = (\sqrt{2} - 1)/2 \approx 0.207$. La tranche a donc pour aire
\[
2\int_0^{y_0} \Bigl(\tfrac14 - y^2 - y\Bigr) dy = 2\Bigl(\frac{y_0}{4} - \frac{y_0^3}{3} - \frac{y_0^2}{2}\Bigr).
\]
En multipliant par 4, on obtient $(4\sqrt{2} - 5)/3$. Comme le montre la \cref{fig:iv1:parabola}, la région est un carré aux côtés bombés. Son aire se situe entre celle du disque de rayon $\tfrac14$, qui vaut environ $0.196$, et celle du disque de rayon $\tfrac12$.

\begin{omfigure}
\begin{tikzpicture}[x=4cm,y=4cm]
\draw (-0.5,-0.5) rectangle (0.5,0.5);
\draw[gray,dashed] (-0.5,-0.5) -- (0.5,0.5);
\draw[gray,dashed] (-0.5,0.5) -- (0.5,-0.5);
\fill[omDef!25,draw=omDef,thick]
 plot[domain=-0.2071:0.2071,samples=30,smooth] ({0.25-\x*\x},{\x})
 -- plot[domain=-0.2071:0.2071,samples=30,smooth] ({-\x},{0.25-\x*\x})
 -- plot[domain=-0.2071:0.2071,samples=30,smooth] ({-(0.25-\x*\x)},{-\x})
 -- plot[domain=-0.2071:0.2071,samples=30,smooth] ({\x},{-(0.25-\x*\x)})
 -- cycle;
\fill (0,0) circle (0.6pt);
\node[font=\footnotesize,below=1pt] at (0,0) {centre};
\end{tikzpicture}

\omcaption{Les points du carré unité plus proches du centre que du bord. Chacun des quatre arcs est une parabole, de foyer le centre et de directrice un côté. Les diagonales en pointillés découpent le carré en quatre tranches.}\label{fig:iv1:parabola}
\end{omfigure}

\iqwrong{Prendre un disque de rayon $\tfrac14$, ce qui donne $\pi/16 \approx 0.196$. Cela n'est juste que le long des axes. En direction d'un coin, le bord est plus éloigné, si bien que la vraie région dépasse le disque, jusqu'à une distance d'environ $0.293$ du centre.}
\end{solution}

\begin{solution}{iq:iv1:continuous-and-geometric-probability:23}
La fraction peinte vaut $7/18 \approx 0.389$.

Pour de nombreux oiseaux, les écarts entre voisins se comportent comme des exponentielles indépendantes\index{approximation de Poisson}. On prend l'écart qui contient un point aléatoire du fil. Il est biaisé par la taille\index{echantillonnage@échantillonnage biaisé par la taille}, donc, avec un écart moyen égal à 1, sa longueur $L$ a pour densité $Le^{-L}$.

L'écart est peint dès qu'au moins l'un de ses deux oiseaux d'extrémité le considère comme son côté le plus court. L'oiseau de gauche ne le peint pas si son autre écart est plus court que $L$, ce qui arrive avec probabilité $1 - e^{-L}$. Les deux autres écarts étant indépendants, l'écart reste non peint avec probabilité $(1 - e^{-L})^2$. En moyennant~:
\[
E[(1 - e^{-L})^2] = \int_0^\infty Le^{-L}\bigl(1 - 2e^{-L} + e^{-2L}\bigr) dL = 1 - \tfrac24 + \tfrac19 = \tfrac{11}{18}.
\]
Ainsi, $11/18$ reste non peint et $7/18$ est peint, ce qu'une simulation d'oiseaux posés sur un fil confirme.

\iqwrong{Affirmer une moitié ou plus, au motif que chaque oiseau peint quelque chose. Les écarts courts sont peints, parfois par les deux extrémités, et les longs tendent à ne pas l'être. Or les longs écarts occupent la plus grande part du fil.}
\iqfollow{Quelle proportion des écarts, en nombre, est peinte~?}{Sans le biais de taille, la densité de $L$ est $e^{-L}$, et la probabilité de rester non peint a pour moyenne $1 - 2 \cdot \tfrac12 + \tfrac13 = \tfrac13$, donc $2/3$ des écarts sont peints. Ce sont les courts, ce qui explique qu'ils n'occupent que $7/18$ de la longueur.}
\end{solution}

\enlargethispage{-3\baselineskip}%
\begin{solution}{iq:iv1:continuous-and-geometric-probability:24}
La probabilité vaut $1/8$.

On tire les points de la façon suivante~: on choisit d'abord quatre droites aléatoires passant par le centre, puis on choisit l'un des deux points de chaque droite à pile ou face\index{théorème de Wendel}\index{argument de symétrisation}. Cela revient à choisir quatre points uniformes indépendants. Il y a 16 choix de signes, et il suffit de compter combien d'entre eux donnent un tétraèdre qui entoure le centre.

On fixe le point $p_4$ sur la quatrième droite. Les trois autres droites donnent chacune deux choix, et ceux-ci forment les 8 triangles sphériques dans lesquels les trois grands cercles (passant par les couples de droites) découpent la sphère. Les huit triangles recouvrent la sphère sans chevauchement, à part les arêtes. Le tétraèdre de sommets $\pm p_1, \pm p_2, \pm p_3, p_4$ contient le centre exactement lorsque $-p_4$ se trouve dans le triangle $\pm p_1, \pm p_2, \pm p_3$. Exactement un des huit triangles contient $-p_4$. La probabilité vaut donc $1/8$, quel que soit $p_4$.

Le même raisonnement sur un cercle donne $1/4$, et plus généralement, en dimension $d$, la probabilité que $d + 1$ points entourent le centre vaut $2^{-d}$. La somme de Wendel étend ce résultat à davantage de points.

\iqwrong{Essayer d'intégrer sur les angles solides. L'astuce supprime tout le calcul. C'est la solution que le recruteur espère voir apparaître, peut-être après le cas bidimensionnel.}
\iqfollow{Et les quatre dans un même hémisphère~?}{Le complémentaire~: $7/8$. La somme de Wendel pour $d = 3$ et $n = 4$ vaut $2^{-3}\bigl(\binom30 + \binom31 + \binom32\bigr) = 7/8$.}
\end{solution}

\begin{solution}{iq:iv1:continuous-and-geometric-probability:25}
Elle croise $4/\pi \approx 1.273$ droites en moyenne, et elle en croise au moins une avec probabilité $3/\pi \approx 0.955$.

L'espérance s'obtient par linéarité\index{linéarité de l'espérance}. La grille est formée de deux familles de droites parallèles, et le résultat de Buffon\index{aiguille de Buffon} donne $2/\pi$ croisements pour chacune, soit $4/\pi$ au total.

Pour la probabilité d'en croiser au moins une, on fixe l'angle $\theta \in (0, \pi/2)$ à l'une des directions~; les autres angles se traitent de la même façon par symétrie. Les positions horizontale et verticale étant indépendantes, l'aiguille évite les deux familles avec probabilité $(1 - \cos\theta)(1 - \sin\theta)$. En moyennant sur $\theta$ de densité $2/\pi$~:
\[
\frac{2}{\pi} \int_0^{\pi/2} (1 - \cos\theta)(1 - \sin\theta) d\theta = \frac{2}{\pi}\Bigl(\frac{\pi}{2} - 1 - 1 + \frac12\Bigr) = 1 - \frac{3}{\pi}.
\]
Le principe d'inclusion--exclusion donne le même résultat~: chaque famille est croisée avec probabilité $2/\pi$, et les deux avec probabilité $\frac{2}{\pi}\int_0^{\pi/2}\cos\theta\sin\theta\, d\theta = 1/\pi$, donc la probabilité d'en croiser au moins une vaut $4/\pi - 1/\pi$.

Une aiguille de cette longueur croise presque à chaque fois une droite, mais elle ne croise des droites des deux directions qu'environ un tiers du temps.

\iqwrong{Utiliser $2 \cdot 2/\pi - (2/\pi)^2$, en traitant les deux familles comme indépendantes. Les positions sont indépendantes, mais l'angle est commun~: une aiguille proche d'une direction croise moins cette famille, et davantage l'autre.}
\iqwhy{Moyenner sur un facteur commun caché crée une dépendance, même lorsque chaque composante est indépendante conditionnellement à lui. Considérer les facteurs communs cachés comme la première chose à vérifier dans tout calcul de probabilité d'événements conjoints.}
\end{solution}
